\section{Python discovery: choosing a computation that preserves the question}
\label{exec13:python-discovery}

A funding analyst can obtain three apparently plausible answers from the same
table: a partially directed graph, a matrix of regression coefficients, and a
distribution over neural structural models. Their common Python interface does
not make these answers interchangeable. The first can preserve directions that
the data have not resolved. The second can describe one selected ordering. The
third can express uncertainty conditional on a generative model, its prior and
an approximation to its posterior. A product that converts all three into a
single directed network would discard precisely the distinctions needed to
decide whether changing a quoted rate is supported by evidence.

The comparison below examines nine implementation families at identified source
revisions preceding the 7 October 2026 cutoff: causal-learn, pgmpy, LiNGAM,
NOTEARS, DAGMA, gCastle, CausalNex, Causica and the CausalDiscoveryToolbox (CDT).
It develops selected static discovery paths, rather than attributing every
algorithm advertised by a package to the code inspected. These are proposed
reuse choices supported by primary implementation and documentation. The
financial and mathematical examples are reproducible project calculations;
they are not fitted outputs from these engines. Broader Python search coverage,
some delegated implementations and empirical comparisons remain open.

\subsection{The funding calculation establishes the matrix contract}

Consider three variables measured over a fixed decision horizon. Let \(F\) be a
change in a quoted funding rate, in percentage points, \(R\) the change in
retained funding, in millions of US dollars, and \(C\) the corresponding change
in contribution margin, in the same monetary units. A hypothetical structural
model is
\begin{equation}
 F=\varepsilon_F,\qquad R=4F+\varepsilon_R,\qquad
 C=-F+\tfrac12R+\varepsilon_C.
 \label{exec13:pydisc-funding}
\end{equation}
The disturbances are centered and mutually independent, and intervention on
\(F\) replaces its first equation. These assumptions define the illustration;
the coefficients are not estimated bank responses. The direct margin cost of
raising the rate is offset by the margin earned on retained funding. A rate
increase of \(0.25\) percentage points therefore changes expected retention by
USD~1 million and expected margin by USD~0.25 million.

For a data matrix whose columns have order \((F,R,C)\), write the row convention
as \(X=XW+E\), with \(W_{ij}\) the coefficient on the edge \(i\to j\). The
corresponding column-vector structural convention is \(x=Bx+e\), where
\(B=W^\top\). In this example,
\begin{equation}
 W=\begin{pmatrix}0&4&-1\\0&0&1/2\\0&0&0\end{pmatrix},\qquad
 T=I+W+W^2=(I-W)^{-1}
  =\begin{pmatrix}1&4&1\\0&1&1/2\\0&0&1\end{pmatrix}.
\end{equation}
Thus \(T_{FC}=1=-1+4(1/2)\) is the total margin response per percentage point.
The matrix identity holds because \(W^3=0\). More generally, a finite DAG with
\(d\) nodes has a nilpotent weighted adjacency matrix and a finite path sum
\(I+W+\cdots+W^{d-1}\). This algebra concerns a supplied linear structural
model; it is not a discovery or identification theorem.

The reference NOTEARS and DAGMA linear losses predict columns through \(XW\).
The inspected LiNGAM adjacency routine instead assigns coefficients to
\texttt{B[target,predictors]} and its path helper multiplies entries in that
convention \citep{exec13_pydisc_notears,exec13_pydisc_dagma,exec13_pydisc_lingam}.
An adapter that interprets LiNGAM's matrix as \(W\) reports zero for the
\(F\to C\) response in this example and invents the reversed path response.
Both matrices are acyclic, both have identical dimensions, and both pass a
check that merely counts edges. Node order, direction convention and units
must therefore be carried as data, verified before any intervention query.

Endpoint matrices require a separate decoder. In causal-learn's PC/GES return
contract, entries \((G_{ij},G_{ji})=(-1,1)\) denote \(i\to j\); a pair of
minus ones denotes an undirected edge. These entries are endpoint marks, not
negative and positive causal coefficients \citep{exec13_pydisc_causallearn}.
Treating them as weights would invent both magnitudes and interventions. The
proposed common result must distinguish a weighted DAG, a partially directed
graph and a posterior over structural models before it exposes any common
plotting or effect-evaluation interface.

\subsection{A finite landscape of complementary mechanisms}

Table~\ref{tab:exec13-pydisc} records a proposed role for each inspected family.
The recommendations depend on the scientific target, not on a single numerical
ranking. A package can be useful as a comparison engine without being the
default dependency, and a method can be worth retaining even when it cannot
resolve the bank's current ambiguity.

\begingroup\small
\begin{longtable}{P{0.16\textwidth}P{0.36\textwidth}P{0.40\textwidth}}
\caption{Selected Python discovery paths and proposed reuse. These are source
comparisons, not measured engine rankings.}\label{tab:exec13-pydisc}\\
\toprule Family & Inspected mechanism & Proposed role and condition \\
\midrule\endfirsthead
\toprule Family & Inspected mechanism & Proposed role and condition \\
\midrule\endhead
causal-learn & PC endpoint graph; GES with selectable scores & Integrate selected
paths with explicit CI, score, collider and graph-type contracts. \\
pgmpy & PC; equivalence search with forward, backward and turning phases &
Integrate the pinned graph API; request partially directed output explicitly. \\
LiNGAM & DirectLiNGAM order selection; coefficient and total-path helpers &
Integrate within a declared linear, non-Gaussian, causally sufficient model;
preserve the target-row convention. \\
NOTEARS & Linear loss, exponential acyclicity and sparsity penalty & Retain
the compact reference as a numerical baseline before considering a port. \\
DAGMA & Linear loss with log-determinant acyclicity and domain control & Retain
a separate baseline; check domain, final topology, input mutation and constraints. \\
gCastle & Selected PC and polynomial-constraint NOTEARS paths & Retain a
matched alternative with method-specific semantics and thresholds. \\
CausalNex & StructureModel and NOTEARS variants with tabu constraints & Reject
as a default new dependency; retain conditional archival comparison. \\
Causica & Learned structural functions, noise and graph distribution & Retain
an experimental generative-model option; evaluate effects conditionally. \\
CDT & Selected PC/GES wrappers invoking R packages & Defer default adoption
and porting until delegated versions, ordering and terms are resolved. \\
\bottomrule
\end{longtable}
\endgroup

The selected PC/GES sources support the first two rows
\citep{exec13_pydisc_causallearn,exec13_pydisc_pgmpy}.
The functional and continuous comparison draws on the LiNGAM and NOTEARS
implementations \citep{exec13_pydisc_lingam,exec13_pydisc_notears}, together
with DAGMA and gCastle \citep{exec13_pydisc_dagma,exec13_pydisc_gcastle}.
The remaining rows draw on the selected CausalNex, Causica and CDT paths
\citep{exec13_pydisc_causalnex,exec13_pydisc_causica,exec13_pydisc_cdt}.
\texttt{castle} is the import namespace of gCastle; it is not an additional
independent implementation. The repository \texttt{cdt15/lingam} and CDT are
distinct projects despite their similar names. Repository names and import
aliases should be resolved before counting competitors or estimating the
cost of integration.

\subsection{Constraint and score methods preserve different qualifications}

The selected causal-learn PC interface exposes the significance threshold,
conditional-independence test, stable skeleton option, collider rule, conflict
priority, background knowledge and conditioning-set bound. It also has a
separate missing-value route. These arguments control substantive choices in
the graph construction. They should appear in the result's explanation whenever
they affect the decision, rather than disappearing into an undifferentiated
``PC'' label \citep{exec13_pydisc_causallearn}.

For its ordinary Fisher test, the source computes a partial correlation from
an inverse correlation submatrix, then forms
\begin{equation}
 Z=\sqrt{n-|S|-3}\,\bigl|\operatorname{atanh}(\widehat r_{ij\mid S})\bigr|,
 \qquad p=2\{1-\Phi(Z)\}.
\end{equation}
Here \(S\) is the conditioning set, \(n\) the row count and \(\Phi\) the
standard normal distribution function. The Gaussian independent-row sampling
model motivates this reference approximation; correlation zero does not imply
conditional independence for an arbitrary non-Gaussian distribution. The code
raises an error for a singular correlation submatrix and clips a computed
correlation at the numerical boundary. Neither clipping nor returning a
number establishes adequate information. A sensible adapter must reject
nonpositive \(n-|S|-3\), constant columns and singular designs with an
explanation, not translate failure into an independence finding.

Repeated observations from one account also challenge this reference law. If a
correlation of \(0.1\) is supplied with 100 independent units, the displayed
statistic is approximately \(0.988\). Substituting 1,000 duplicated rows gives
approximately \(3.168\), although the duplicated data contain no new
independent information. The example compares the formula under two declared
row counts; it does not propose a universal effective-sample-size correction
\citep{exec13_pydisc_examples}. A panel or time-series request must use a
sampling argument appropriate to that dependence, or be declined by this
static interface.

The inspected pgmpy PC path defaults to a partially directed graph and limits
conditioning sets to five variables unless configured otherwise. Its optional
DAG is one completion of the learned partial graph. Its documented background
knowledge includes required and forbidden directions, while collider conflicts
can still generate a warning. Consequently the proposed adapter should check
the returned graph against declared constraints instead of treating the mere
presence of an expert-knowledge argument as proof that every input condition
was satisfied \citep{exec13_pydisc_pgmpy}. A five-variable search limit is a
computational scope choice, not an assertion that larger separators cannot
matter in the bank's model.

Score methods make a different tradeoff. causal-learn GES accepts several
local scores, including Gaussian BIC, discrete BDeu and regression-based
alternatives. Its inspected BIC routine computes residual variance and can use
a pseudoinverse or a small positive replacement for a nonpositive numerical
variance. pgmpy's inspected GES performs forward, backward and turning moves,
uses a minimum improvement threshold, and converts its final equivalence graph
to a DAG only if requested
\citep{exec13_pydisc_causallearn,exec13_pydisc_pgmpy}. A larger score in one
engine cannot be compared directly with a number from another engine without
matching sign, constants, likelihood, penalties and preprocessing. The
proposal should compare graph/effect consequences under matched assumptions,
not select an engine because its default score is numerically larger.

\subsection{Functional restrictions can orient graphs, but must be stated}

LiNGAM is attractive when the banking question justifies a linear acyclic model
with independent non-Gaussian disturbances and no unobserved confounders. Those
are explicit premises in the original work, including independent repeated
observations from the same model \citep{exec13_pydisc_shimizu2006}. The original
paper describes an ICA route; the current DirectLiNGAM implementation inspected
here selects an ordering through residual comparisons, then estimates
coefficients using that order. The original result should not be quoted as
if it had independently verified every modern option and integration.

The selected helper standardizes variables for an adaptive-Lasso pruning step,
uses a BIC criterion, and refits the retained predictors in original units.
The separate total-effect helper sums directed path products and checks that
variables along those paths satisfy its continuous-variable condition
\citep{exec13_pydisc_lingam}. These are valuable implementation details for a
funding adapter: they explain where a sparse edge set comes from and how the
reported monetary slope is recovered. They do not show that bank disturbances
are independent or non-Gaussian. Heavy tails in an observed balance series
alone do not establish the required disturbance model after conditioning on
the structural parents.

The observational ambiguity introduced earlier provides an exact adverse case.
For a standardized pair, consider two centered Gaussian models,
\begin{align}
 \mathcal M_{\to}:&\quad X=U,\quad Y=0.6X+V,
     &\operatorname{var}(U)&=1,&\operatorname{var}(V)&=0.64,\\
 \mathcal M_{\leftarrow}:&\quad Y=\widetilde U,\quad X=0.6Y+\widetilde V,
     &\operatorname{var}(\widetilde U)&=1,&\operatorname{var}(\widetilde V)&=0.64.
\end{align}
Disturbances are independent within each model. Both laws have covariance
\(\left(\begin{smallmatrix}1&0.6\\0.6&1\end{smallmatrix}\right)\), so the
complete observational distributions coincide, not just a selected correlation.
The effect of intervening on \(X\) is nevertheless \(0.6\) in the first model
and zero in the second. Within this declared two-model class the effect set is
\(\{0,0.6\}\); the interval \([0,0.6]\) is its convex hull, not an assertion
that every intermediate effect belongs to the two-model class.

No algorithm using observations alone can distinguish these two Gaussian laws
reliably in both worlds. A non-Gaussian identifying premise, a justified
temporal direction, or an intervention changes the available information.
An optimizer's ability to return one orientation does not. The proposed
DirectLiNGAM integration therefore retains this example as a refusal or
assumption-failure case, rather than interpreting a selected order as empirical
proof of its own premises.

\subsection{Continuous acyclicity is a numerical condition with a domain}

Continuous methods replace a combinatorial acyclicity constraint by a smooth
one. Let \(S=W\circ W\), the entrywise square of \(W\). NOTEARS uses
\begin{equation}
 h_{\exp}(W)=\operatorname{tr}(\exp S)-d
             =\sum_{k\geq1}\frac{\operatorname{tr}(S^k)}{k!}.
\end{equation}
Because \(S\) is nonnegative, its diagonal powers accumulate nonnegative
closed-walk products. A directed cycle creates a positive term; in a DAG all
such terms vanish. This is the exact-zero characterization developed in
\citet{notears2018}. Its numerical implementation uses tolerances, iteration
limits and thresholding, each of which must be distinguished from the exact
statement.

The source inspected for gCastle's class named \texttt{Notears} instead uses
the polynomial expression
\begin{equation}
 h_{\rm poly}(W)=\operatorname{tr}\{(I+S/d)^d\}-d.
\end{equation}
The reference NOTEARS linear implementation uses a matrix exponential. The
two expressions share an exact-zero acyclicity property but have different
values, derivatives and numerical scaling away from zero. gCastle also keeps
the weighted estimate separately from a binary matrix formed with the strict
comparison \(|W_{ij}|>\tau\), whereas the reference removes weights satisfying
\(|W_{ij}|<\tau\). At equality, the reference retains the edge and gCastle's
binary output does not \citep{exec13_pydisc_gcastle,exec13_pydisc_notears}.
A benchmark that silently compares one weighted matrix with another binary
matrix can mistake a representation choice for a scientific result.

DAGMA uses the log-determinant characterization
\begin{equation}
 h_{\log,s}(W)=-\log\det(sI-S)+d\log s,
 \qquad s>\rho(S),
\end{equation}
where \(\rho(S)\) is the spectral radius. The domain ensures that \(sI-S\) is
a nonsingular M-matrix. The domain condition is part of the theorem, not an
optional numerical detail \citep{exec13_pydisc_dagma2022}. For example, within
this domain one can expand the log determinant as
\(\sum_{k\geq1}\operatorname{tr}(S^k)/(k s^k)\); its nonnegative cycle terms
again explain the zero condition. Computing a log absolute determinant outside
the domain does not justify the same conclusion.

The inspected linear DAGMA implementation uses inverse-matrix checks and
backtracking in its optimization. Its documentation explicitly allows the
possibility that the finite output is cyclic and states that requested included
edges are not guaranteed to survive. It also centers the supplied array in
place and saves a final residual before thresholding
\citep{exec13_pydisc_dagma}. These findings motivate copying the input,
retaining the full optimization configuration, and recomputing topology and
residual diagnostics on the actual returned graph. They establish a proposed
adapter requirement; they do not establish a measured failure rate.

A small exact counterexample illustrates why a tolerance cannot replace a
topological check. Set
\(W=\left(\begin{smallmatrix}0&\epsilon\\\epsilon&0\end{smallmatrix}\right)\)
with \(\epsilon=0.005\). This is a two-cycle. Nevertheless,
\begin{align}
 h_{\exp}(W)&=2\{\cosh(\epsilon^2)-1\}\simeq6.25\times10^{-10},\\
 h_{\rm poly}(W)&=\epsilon^4/2=3.125\times10^{-10},\\
 h_{\log,1}(W)&=-\log(1-\epsilon^4)\simeq6.25\times10^{-10}.
\end{align}
All are below \(10^{-8}\). A threshold of \(0.001\) retains both edges. The
reference default threshold of \(0.3\) removes both; this example is therefore
not a claim that a default solver necessarily returns this cycle. It shows
why the numerical residual, thresholded support and scientific effect are
three distinct objects \citep{exec13_pydisc_examples}.

\subsection{Units and centering can change the selected orientation}

In the preceding two-model Gaussian example, the population least-squares
loss \(\tfrac12\sum_j\operatorname{var}(\varepsilon_j)\) is \(0.82\) in
both orientations, and the coefficient \(\ell_1\) penalty is \(0.6\) in both.
Both graphs are acyclic. Thus neither perfect acyclicity nor this matched
population objective distinguishes the two causal effects. Scores that add
other model restrictions may distinguish them, but those restrictions must be
declared and defended.

Now express \(Y\) in units ten times as large numerically, \(Y'=10Y\), while
leaving \(X\) unchanged. The forward coefficient becomes \(6\), the reverse
coefficient becomes \(0.06\), and their residual-variance sums become \(65\)
and \(100.64\), respectively. The corresponding least-squares losses are
\(32.5\) and \(50.32\). The physical observational ambiguity is unchanged,
but the unweighted loss and sparsity tradeoff have changed
\citep{exec13_pydisc_examples}. This matters when a funding dataset mixes
percentage points, currency, duration and account counts. Standardization is a
modeling decision that must be recorded and reversed when reporting monetary
responses; it does not independently supply causal identification.

CausalNex's selected unpenalized NOTEARS path applies its squared-error
criterion to the supplied numeric array without the centering performed in the
reference NOTEARS least-squares path. It provides tabu edges and prohibited
parent/child nodes as optimization bounds, and its StructureModel records edge
origins. The same graph object has an explicit operation that repeatedly drops
the smallest absolute-weight edge until a DAG remains, with a warning about
different units \citep{exec13_pydisc_causalnex}. Dropping an edge changes the
model used by an intervention query. It must not be treated as invisible file
cleanup, and the thresholded model should be assessed on its own assumptions.

The pinned CausalNex package declares Python versions from 3.8 to below 3.11.
That declaration does not prove installation failure in every newer
environment, but it is sufficient to reject an unqualified promise that the
current common Python environment can adopt it directly. The maintenance
review in Section~\ref{exec14:causalnex} further rules it out as a default new
dependency under the inspected end-of-life notice. An isolated process remains
a possible archival comparison when a concrete reproducibility need justifies
an accountable maintainer and the additional environment. For the simpler
linear reference comparison, wrapping NOTEARS directly is a substantive
alternative that avoids introducing that framework dependency merely to
obtain a DAG matrix.

\subsection{Causica connects discovery and effects through a learned model}

Causica's inspected DECI path jointly learns a distribution over graphs,
functional relationships and noise. Its training step samples a relaxed graph,
evaluates a data log density, incorporates the graph distribution's entropy
and a graph prior, and adds an acyclicity penalty. Suppressing implementation
scaling details, the statistical part has the form
\begin{equation}
 -\mathbb E_{q_\phi(G)}\log p_\theta(X\mid G)
 +\operatorname{KL}\{q_\phi(G)\,\|\,p(G)\}.
\end{equation}
Here \(q_\phi\) is an approximating graph distribution, \(\theta\) parameterizes
the structural functions and noise, and \(p(G)\) contains graph preferences.
The actual module uses minibatches, relaxed samples and scaled penalties;
this expression explains their statistical roles, not an assertion of exact
posterior computation. The inspected prior includes sparsity and an optional
weighted discrepancy from a supplied expert graph
\citep{exec13_pydisc_causica}.

For an assumed structural model, simulation after replacing an intervention
equation can estimate a model-implied ATE. The Causica documentation describes
this workflow, while the inspected test methods separately evaluate graph
reconstruction, observational likelihood, interventional likelihood and effect
errors when appropriate reference data are supplied. These distinctions are
useful for the proposed product. A model may predict observed balances well
while choosing the wrong intervention direction in an observationally ambiguous
case; graph uncertainty under \(q_\phi\) is conditional on the model family,
prior and approximation. It is not automatically a confidence set with
frequentist coverage for a bank effect.

In particular, an expert graph supplied by a language model should retain its
origin and uncertainty. A strong prior favoring a rate-to-balance direction can
concentrate a fitted graph distribution without providing new assignment
evidence. A suitable prototype would compare the same task under an empty
prior, a justified restricted prior and a deliberately wrong prior, preserving
the resulting change in intervention conclusions. Such a comparison is a
proposed experiment, not a result of this source review.

The exact revision also matters for capability claims. Its README states that
Rhino and DDECI were removed from the refactored codebase and points to an older
release. Their historical association with Causica cannot be used to claim
that the inspected DECI path currently implements general temporal or
latent-confounded inference. A missingness or latent-variable claim needs the
corresponding implementation and assumptions, examined in its own right.

\subsection{A Python wrapper can carry another language's scientific contract}

CDT's selected PC and GES classes illustrate why language should be recorded
at both the public interface and the computational backend. PC checks for
\texttt{pcalg}, \texttt{kpcalg} and \texttt{RCIT}; GES checks for
\texttt{pcalg}. The wrappers construct files, dispatch R templates and read
results into Python graph objects \citep{exec13_pydisc_cdt}. A Python package
pin alone therefore does not freeze the statistical engine, the R environment
or its graph encoding. A direct R adapter is a serious alternative if it
preserves the same semantics with fewer conversion layers.

There are two concrete reasons to defer a default CDT recommendation. First,
its GES documentation labels one option as categorical, while the code maps
the corresponding \texttt{int} key to \texttt{GaussL0penIntScore}. The selected
Python wrapper does not establish a discrete likelihood merely because of
that prose label. Section~\ref{sec:exec28-bp} follows the selected R template:
it constructs the score from the data alone. The separately inspected pcalg
class is Gaussian, with default empty intervention targets. This narrows
the capability question but does not establish which pcalg version a CDT
installation resolves. The wrapper also relabels graph nodes using sorted
column names in a selected path. The later exact permutation example explains
why deliberately unsorted columns require a native conversion test before
banking use. \citep{exec28_bp_cdt}

Second, the repository's saved license text is MIT, while the saved setup
metadata names Apache~2.0. Both observations are retained. Neither is a reason
to invent a single reconciled distribution license or to assume delegated
R packages inherit the wrapper's terms. This is a local adoption question,
not a reason to stop comparisons of unrelated Python implementations. The
other selected primary terms are Apache~2.0 for NOTEARS, DAGMA and the inspected
gCastle root; CausalNex supplies Apache terms and separate trademark text;
Causica, LiNGAM and the pinned pgmpy license supply MIT text. The initial
causal-learn notice came from a mutable-main snapshot. The reuse review in
Section~\ref{exec14:causallearn-decision} also retains the MIT notice at the
selected technical revision. That resolves the root-notice pin; it does not
resolve the complete dependency and distribution profile. These observations
guide dependency review; they do not constitute a complete transitive-license
opinion. \citep{exec14_causallearn_reuse}

\subsection{Selection and missingness change what a discovery table means}

A funding table may contain only customers who were successfully contacted,
retained an account or completed a follow-up. Treating the remaining rows as a
random sample can create edges even before choosing a discovery algorithm.
For an exact generated illustration, let \(F\) and \(N\) be independent
standard normal variables and observe only rows satisfying \(F+N>0\). Define
\(U=(F+N)/\sqrt2\) and \(V=(F-N)/\sqrt2\). Before selection, \(U\) and \(V\)
are independent standard normals. After selection, \(U\) is half-normal with
variance \(1-2/\pi\), while \(V\) still has variance one. It follows that
\begin{equation}
 \operatorname{cov}(F,N\mid F+N>0)=-1/\pi,\qquad
 \operatorname{corr}(F,N\mid F+N>0)=-\frac1{\pi-1}.
\end{equation}
The selected correlation is about \(-0.467\), despite population independence
\citep{exec13_pydisc_examples}. A graph of the selected table can correctly
describe its associations while misleading a population intervention query.

This example gives practical meaning to the distinction between a missing-value
option and an argument for selection recovery. A package-specific route must
state which response mechanism it handles and what observations identify its
correction. Imputation, complete-case analysis and a generative missing-data
model make different commitments. The proposed request should carry the
sampling unit, observation window, selection indicators, and relevant
conditioning events. If these are unavailable, the product can still support
exploratory associations and alternative diagrams, but its intervention
conclusion must retain the unresolved selection assumption.

\subsection{An adapter contract and an economical prototype choice}

The proposed static-discovery boundary first validates the question and data:
units, node order, information timing, row dependence, missingness, allowed
interventions and model family. It then configures one pinned engine and
retains its native result. A second step decodes that result into a typed graph
or structural-model object. Only a subsequent identification step may turn
that object into a causal effect or effect set. For a partially directed
result, selecting one completion is an explicit additional assumption; for a
neural SEM, reporting simulation output is conditional on the selected learned
model and uncertainty representation.

Useful deterministic acceptance cases follow directly from the comparison.
The funding model tests a transpose and a deliberate permutation of columns.
The two Gaussian orientations test whether an interface preserves unresolved
effects. The weak two-cycle tests support-level acyclicity independently of a
floating-point residual. A coefficient exactly on the threshold tests the
inequality convention. A singular CI input tests refusal, while duplicated
rows test whether the declared sampling unit is preserved. These cases check
the proposed boundary; they are not substitutes for a prospectively designed
empirical evaluation of discovery accuracy or uncertainty.

The initial prototype should therefore favor a small, interpretable set of
mechanisms: one PC/GES path that preserves equivalence, one functional path
with explicit non-Gaussian assumptions, and separate NOTEARS/DAGMA numerical
references. gCastle can test the cost of a broader algorithm catalog.
CausalNex remains an archival option under the maintenance conditions in
Section~\ref{exec14:causalnex}. CDT remains conditional on a consumer need
that justifies its additional dependency and conversion burden. Causica supplies a distinct
generative-model experiment, whose contribution should be judged on matched
intervention questions and negative cases rather than an attractive graph
alone. Porting a numerical core becomes worthwhile only when measured
deployment constraints or performance needs exceed the cost of validating a
second implementation.

These choices make the proposal concrete without assigning a fictional bank
ROI. Their economic comparison must include environment maintenance, source
and assumption review, translation errors, failed requests and the cost of
resolving consequential graph ambiguity. The known routes are exact interface
tests, scoped primary-source completion and prospectively designed public or
generated comparisons. The unresolved scientific question is whether the
information available for a particular bank decision supports the required
intervention claim. Adding another optimizer does not answer that question.
\label{exec13:python-discovery-end}
